\section{Model and Computational Methodology}
\label{sec:methodology}

This chapter is the implementation contract for the completed numerical audit.
It fixes the physical variables, solver-coordinate map, stenosis geometry, wall
law, source terms, boundary approximation, principal finite-volume realization,
velocity observation operator, comparison design, and provenance record before
any numerical result is interpreted. Secondary solver surfaces remain appendix
and verification context unless a result table explicitly identifies them as the
active method.

\subsection{Physical and Solver Variables}
\label{subsec:methodology-problem-definition}
\label{subsec:selected-1d-forward-model}
\label{subsec:cross-sectional-variables-governing-equations}
\label{subsubsec:cross-sectional-reduction-selected-1d-model}

The computational problem is a single idealized stenotic vessel represented by a
selected one-dimensional area--flow model. The model assumes a distinguished
axial direction and cross-sectionally averaged pressure and velocity, retaining
distributed pulse-wave and pressure--flow structure while discarding secondary
flow, separation, recirculation, and resolved wall traction. The reduction
follows standard compliant-artery and dimension-reduced modeling arguments
\cite[Sec.~2, pp.~253--257]{FormaggiaEtAl2003OneDimensionalBloodFlow}
\cite[Ch.~3, Secs.~3.2 and~3.6]{KopplHelmig2023DimensionReducedModeling}
\cite[Sec.~3, pp.~24--26]{ShiEtAl2011BloodFlowModelReview}.

\begin{assumption}[Straight single-vessel geometry]
    \label{ass:1d-straight-vessel-geometry}
    The selected 1D model assumes a vessel of length $L$ aligned with the
    $z$-axis and a circular cross-section with radius $R(z,t)$. For
    $0<z<L$, $S(z,t)$ is the interior cross-section used for averaging, not
    a boundary component of the open lumen.
\end{assumption}

\begin{definition}[Section-averaged 1D fields]
    \label{def:section-averaged-1d-fields}
    The retained reduced variables are
    \begin{flalign*}
        \quad A_{\mathrm{phys}}(z,t) &\defined \int_{S(z,t)} 1 \, dS, &&
            \text{physical cross-sectional area}, \\
        \quad Q_{\mathrm{phys}}(z,t) &\defined \int_{S(z,t)} u_z \, dS, &&
            \text{physical volumetric flow rate}, \\
        \quad \bar u(z,t) &\defined
            \frac{Q_{\mathrm{phys}}(z,t)}{A_{\mathrm{phys}}(z,t)}, &&
            \text{mean axial velocity}, \\
        \quad p_{\mathrm g}(z,t) &\defined
            \frac{1}{A_{\mathrm{phys}}(z,t)}
            \int_{S(z,t)} p_{3D,\mathrm g} \, dS, &&
            \text{mean gauge pressure}.
    \end{flalign*}
    Here $p_{3D,\mathrm g}$ is continuum pressure in the wall-law gauge frame.
\end{definition}

\begin{definition}[Physical-to-solver map for the implemented 1D state]
    \label{def:physical-to-solver-map}
    The continuum 1D model uses physical area $A_{\mathrm{phys}}$ and
    physical flow $Q_{\mathrm{phys}}$. The implemented radius-squared solver
    instead stores
    \begin{flalign*}
        \quad
        a=R^2,\qquad q=\frac{Q_{\mathrm{phys}}}{\pi}.
        &&
    \end{flalign*}
    Thus
    \begin{flalign*}
        \quad
        A_{\mathrm{phys}}=\pi a,\qquad
        Q_{\mathrm{phys}}=\pi q,\qquad
        \bar u=\frac{Q_{\mathrm{phys}}}{A_{\mathrm{phys}}}=\frac{q}{a}.
        &&
    \end{flalign*}
    With $K=Eh/(1-\sigma^2)$ and the Canic conservative-form convention
    $R_0^*=R_{\max}$, the selected evolution wall law gives the solver elastic
    flux contribution recorded in Appendix~\ref{app:num-current-solver-operator}:
    \begin{flalign*}
        \quad
        \Psi_h(a)=\frac{K}{3\rho R_{\max}^2}a^{3/2},
        \qquad
        \pd_a\Psi_h(a)=\frac{K}{2\rho R_{\max}^2}\sqrt a.
        &&
    \end{flalign*}
    This map is the convention used by the Section~\ref{sec:resolved-3d-comparison}
    velocity and flow comparisons.
\end{definition}

\subsection{Geometry and Wall Law}
\label{subsec:stenosis-geometry-implemented-coordinates}
\label{subsec:selected-wall-law-closure-model}

\begin{definition}[Stenosis reference geometry]
    \label{def:1d-stenosis-reference-geometry}
    A stenosis is encoded through the reference radius or reference area. Let
    $R_{\mathrm{base}}(z)$ be the nominal healthy radius and $s(z)\in[0,1)$
    a fractional radius-reduction profile. Then
    \begin{flalign*}
        \quad
        R_0(z) &= R_{\mathrm{base}}(z)\big(1-s(z)\big),
        &
        A_0(z) &= \pi R_0(z)^2.
        \numberthis\label{eq:1d-stenosis-area} &&
    \end{flalign*}
    The baseline idealized stenotic-vessel profile is the smooth asymmetric
    kernel
    \begin{flalign*}
        \quad
        g(z) &= z-z_s+\kappa
        \exp\left(-\frac{(z-z_a)^2}{2\sigma_a^2}\right),
        &
        \psi(z) &= \exp\left(-\lambda g(z)^4\right),
        \numberthis\label{eq:1d-stenosis-kernel} && \\
        \quad
        s(z) &= s_{\max}\psi(z), \qquad 0\leq s_{\max}<1.
        \numberthis\label{eq:1d-stenosis-profile} &&
    \end{flalign*}
    In the default centimeter-scale simulation record, $z_a=2.5\,\mathrm{cm}$,
    $z_s=3.4\,\mathrm{cm}$, $\sigma_a=1\,\mathrm{cm}$,
    $\kappa=0.95\,\mathrm{cm}$, and
    $\lambda=50\,\mathrm{cm}^{-4}$.
    This geometry is the baseline idealized simulation geometry in this report,
    not patient-specific morphology.
\end{definition}

\begin{remark}[Stenosis-profile and variable-radius boundary]
    \label{rmk:1d-stenosis-profile-scope}
    The profile in Eq.~\ref{eq:1d-stenosis-profile} is $C^\infty$ and
    rapidly localized rather than compactly supported. This smooth idealized
    geometry is the baseline for the reduced simulations in this report and is
    also useful for fixed-wall 3D studies: the same analytic radius profile
    supplies controlled derivatives, wall normals, severity variation, and
    repeatable mesh-refinement families without introducing segmentation noise.
    Canic et al.\ show that standard 1D variable-geometry models may miss
    stenotic pressure and velocity trends unless additional variable-radius
    terms are included \cite[Abstract; Secs.~1--2]{CanicEtAl2024Extended1DStenosis}.
    Those corrections are included only by the \texttt{canic-extended-1d}
    forward-model descriptor; the classical parabolic-profile baseline
    currently recorded by the historical token \texttt{classical-1d-no-slip}
    retains the same selected wall law but omits the Canic variable-radius
    corrections.
\end{remark}

\begin{definition}[Classical parabolic-profile 1D baseline]
    \label{def:classical-no-slip-1d-baseline}
    The classical baseline used as a model-family member is a
    parabolic-profile reduced model. A no-slip wall trace is compatible with a
    parabolic profile for fully developed Newtonian flow in a straight circular
    tube, but no-slip alone does not imply that profile in an unsteady stenotic
    vessel. The parabolic profile is therefore an independent cross-sectional
    closure with $\alpha_{\mathcal P}=4/3$ and $g_{\mathcal P}=4$. In the
    implementation manifests this member is still recorded under the historical
    token \texttt{classical-1d-no-slip}; in the mathematical discussion it is
    the classical parabolic-profile baseline. It keeps the same reduced state
    variables and selected Canic-derived pressure--area wall law as the
    extended 1D framework, but omits the Canic effective-alpha variable-radius
    correction. The descriptor \texttt{canic-extended-1d} is the historical
    manifest token for the current $R_{\max}$-normalized Canic-derived
    extended-model member.
\end{definition}

\begin{definition}[Selected Canic--Koiter pressure--area wall law]
    \label{def:1d-elastic-wall-law}
    The selected elastic wall closure is the Canic--Koiter thin-wall
    pressure--area law. In conventional area notation it may be written
    \begin{flalign*}
        \quad p_{\mathrm g}(A,z,t)
        =
        p_{\mathrm{ext}}(z,t)
        + \frac{\beta(z)}{A_0(z)}
            \left(\sqrt{A(z,t)}-\sqrt{A_0(z)}\right).
        \numberthis\label{eq:1d-wall-law} &&
    \end{flalign*}
    Here $A_0(z)$ is the reference area, $p_{\mathrm{ext}}$ is the external
    pressure, and $\beta(z)$ is a wall-stiffness parameter. In the
    radius-squared coordinate $a=R^2$ used by the implemented finite-volume
    evolution, the selected continuous source-balanced member uses the Canic
    conservative-form reference radius $R_0^*=R_{\max}$:
    \begin{flalign*}
        \quad
        p_{1,\mathrm g}(a,z,t)
        =
        p_{\mathrm{ext}}(z,t)
        +
        \frac{K}{R_{\max}^2}
        \left(\sqrt a-R_0(z)\right),
        \qquad
        K=\frac{Eh}{1-\sigma^2}. &&
    \end{flalign*}
    With $A=\pi a$ and $A_0=\pi R_0^2$, the conventional law in
    Eq.~\ref{eq:1d-wall-law} reduces to this selected member only when
    \begin{flalign*}
        \quad
        \beta(z)=\sqrt{\pi}\,K\,\frac{R_0(z)^2}{R_{\max}^2}. &&
    \end{flalign*}
    Thus the implemented evolution uses a constant reference-radius
    denominator, not the local $R_0^{-2}$ denominator. In this report
    $R_{\max}$ is the healthy reference-radius parameter used by the solver;
    for the default stenosis profile it equals $\max_zR_0(z)$. The selected
    comparison runs use the cgs numerical-record values
    $\rho=1.055\,\mathrm{g}/\mathrm{cm}^3$,
    $E=5.02\times10^6\,\mathrm{dyn}/\mathrm{cm}^2$,
    $h=0.06\,\mathrm{cm}$, $\sigma=0.5$,
    $K=4.016\times10^5\,\mathrm{dyn}/\mathrm{cm}$,
    $R_{\max}=0.18\,\mathrm{cm}$, and
    $p_{\mathrm{ext}}\equiv0$.
    This is the wall law paired with the implemented elastic potential and
    wall/geometry source below. Legacy pressure-output helpers in the
    implementation still report a local-denominator pressure diagnostic through
    \path{pressure(result, params)}; benchmark pressure extrema and
    pressure-diagnostic bars refer to that diagnostic, not to the selected
    evolution wall law.
\end{definition}

\subsection{Implemented $R_{\max}$-Normalized Balance Law}
\label{subsec:conservative-form-source-terms}

\begin{definition}[Closed 1D area-flow model]
    \label{def:closed-1d-area-flow-model}
    For $\mathcal W=\mathcal W_{\mathrm{elastic1D}}$ and an incompressible
    fluid with density $\rho$, the working compliant single-vessel model is
    written with the composite axial pressure derivative
    \begin{flalign*}
        \quad
        \frac{d}{dz}p_{\mathrm g}(A(z,t),z,t)
        =
        \pd_A p_{\mathrm g}(A,z,t) \pd_z A
        +
        \left.\pd_z p_{\mathrm g}(A,z,t)\right|_A. &&
    \end{flalign*}
    The reduced balance law is
    \begin{flalign*}
        \quad \pd_t A
        + \pd_z Q
        &= 0, \numberthis\label{eq:1d-mass-aq} && \\
        \quad \pd_t Q
        + \pd_z
            \left(\alpha \frac{Q^2}{A}\right)
        + \frac{A}{\rho}\frac{d}{dz}p_{\mathrm g}(A(z,t),z,t)
        &= - C_f(z,A,Q) \frac{Q}{A}.
        \numberthis\label{eq:1d-momentum-aq} &&
    \end{flalign*}
    The parameter $\alpha$ is the momentum-flux correction factor induced by
    the selected cross-sectional velocity profile, $C_f$ is the wall-friction
    closure, and $p_{\mathrm g}(A,z,t)$ is the wall-law gauge pressure closing
    the reduced system.
\end{definition}

\begin{assumption}[Velocity-profile and rheology closure]
    \label{ass:1d-first-run-alpha-friction}
    A closed 1D area-flow model must declare the momentum-flux correction
    factor $\alpha$, the wall-friction closure $C_f$, any rheology descriptor
    $\mathcal R$, and the pressure--area wall closure $\mathcal W$. These
    choices determine how cross-sectional profile effects, wall friction,
    viscosity or effective viscosity, and wall compliance enter the reduced
    equations. Numerical-record details such as solver coordinates,
    shear-rate proxies, and discrete source terms are deferred to
    Appendix~\ref{app:numerical-methods-details}.
\end{assumption}

\begin{definition}[Conservative flux/source notation for the reduced wall law]
    \label{def:1d-conservative-flux-source-notation}
    For a reduced wall closure $\mathcal W$ admitting an elastic potential,
    define
    \begin{flalign*}
        \quad
        \mathbf U
        &=
        (A,Q)^\top, &&
        \mathbf F_{\mathcal W}(\mathbf U,z,t)
        =
        \begin{pmatrix}
            Q\\[2pt]
            \alpha Q^2/A+\Psi_{\mathcal W}(A,z,t)
        \end{pmatrix}, &&\\
        \quad
        \mathbf S_{\mathcal W,\mathcal R}(\mathbf U,z,t)
        &=
        \begin{pmatrix}
            0\\[2pt]
            S_{\mathrm{geom}}^{\mathcal W}(A,z,t)
            +S_f^{\mathcal R}(A,Q,z)
        \end{pmatrix}. &&
    \end{flalign*}
    The potential and sources are tied to the selected closure descriptors by
    \begin{flalign*}
        \quad
        \pd_A \Psi_{\mathcal W}(A,z,t)
        &=
        \frac{A}{\rho} \pd_A p_{\mathrm g}^{\mathcal W}(A,z,t), &&\\
        \quad
        S_{\mathrm{geom}}^{\mathcal W}(A,z,t)
        &=
        \left.\pd_z \Psi_{\mathcal W}(A,z,t)\right|_A
        -
        \frac{A}{\rho}
        \left.\pd_z p_{\mathrm g}^{\mathcal W}(A,z,t)\right|_A,
        &&\\
        \quad
        S_f^{\mathcal R}(A,Q,z)
        &=
        -C_f^{\mathcal R}(z,A,Q)\frac{Q}{A}. &&
    \end{flalign*}
    For $\mathcal W_{\mathrm{elastic1D}}$,
    $p_{\mathrm g}^{\mathcal W}$ is the wall law in
    Definition~\ref{def:1d-elastic-wall-law}. In component form,
    \begin{flalign*}
        \quad \pd_t A + \pd_z Q &= 0,
        \numberthis\label{eq:1d-conservative-mass} && \\
        \quad \pd_t Q
        + \pd_z \left(
            \alpha\frac{Q^2}{A}+\Psi_{\mathcal W}(A,z,t)
        \right)
        &=
        S_{\mathrm{geom}}^{\mathcal W}(A,z,t)
        +S_f^{\mathcal R}(A,Q,z),
        \numberthis\label{eq:1d-conservative-momentum} &&
    \end{flalign*}
    The displayed split records the balance-law objects needed by finite-volume
    schemes, while a computed realization may still use a particular
    well-balanced or source-split discretization
    \cite{LeVeque2002FiniteVolumeMethodsForHyperbolicProblems}
    \cite[Summary; Secs.~2.1--2.3; Secs.~5--6]{BoileauEtAl2015Benchmark}
    \cite{FormaggiaEtAl2007FSICoupling}.
\end{definition}

\begin{definition}[Selected implemented balance law]
    \label{def:selected-implemented-balance-law}
    The selected one-dimensional evolution model advanced by the finite-volume
    solver is stated in solver coordinates:
    \begin{flalign*}
        \quad \pd_t a+\pd_z q &=0, &&\\
        \quad \pd_t q
        +\pd_z\left[
            \alpha_{\mathrm{eff}}(z)\frac{q^2}{a}
            +\Psi_h(a)
        \right]
        &=
        S_{\mathrm{wall}}(a,z)
        +S_{\mathrm{fric}}(a,q,z)
        +S_{\alpha}(a,q,z)
        +S_{p_2}(a,q,\pd_z a,\pd_z q,z). &&
    \end{flalign*}
    The terms are
    \begin{flalign*}
        \quad
        \alpha_{\mathrm{eff}}(z)
        &=
        \alpha_{\mathcal P}
        +\iota_{\mathcal M}\alpha_c(R_0'(z)),
        &
        \alpha_c(R_0')&=-\frac{2}{35}(R_0')^2, &&\\
        \quad
        \Psi_h(a)
        &=
        \frac{K}{3\rho R_{\max}^2}a^{3/2},
        &
        S_{\mathrm{wall}}(a,z)
        &=
        \frac{K}{\rho R_{\max}^2}aR_0'(z), &&\\
        \quad
        S_{\mathrm{fric}}(a,q,z)
        &=
        -2\nu_{\mathrm{eff}}^{\mathcal R}
        g_{\mathcal P}\frac{q}{a},
        &
        S_{\alpha}(a,q,z)
        &=
        \iota_{\mathcal M}\frac{q^2}{a}\alpha_c'(z), &&\\
        \quad
        S_{p_2}(a,q,\pd_z a,\pd_z q,z)
        &=
        -\iota_{\mathcal M}a\,\partial_z^{\mathrm{frz}}p_{2,\mathrm g}(a,q,z).
        &&
    \end{flalign*}
    The pressure correction $p_{2,\mathrm g}$ entering this frozen derivative
    is defined in Section~\ref{subsec:source-to-implementation-variant}.
    The indicator $\iota_{\mathcal M}=1$ gives the
    \texttt{canic-extended-1d} member and $\iota_{\mathcal M}=0$ gives the
    classical parabolic-profile baseline. The symbol
    $\partial_z^{\mathrm{frz}}p_{2,\mathrm g}$ denotes the density-scaled
    derivative of the Canic $p_2$ pressure correction with
    $\nu_{\mathrm{eff}}^{\mathcal R}$ held locally fixed during the derivative,
    \begin{flalign*}
        \quad
        \partial_z^{\mathrm{frz}}p_{2,\mathrm g}
        \defined
        \frac{1}{\rho}
        \left.
        \frac{d}{dz}p_{2,\mathrm g}
        \right|_{\nu_{\mathrm{eff}}^{\mathcal R}\ \mathrm{held\ fixed}}.
        &&
    \end{flalign*}
    The selected model therefore contains a derivative-dependent
    nonconservative correction, realized discretely through $D_z^h a$ and
    $D_z^h q$ in Appendix~\ref{app:num-current-solver-operator}, rather than a
    purely local source $S(U,z)$. Thus the classical baseline is recovered by
    setting
    $\alpha_{\mathrm{eff}}=\alpha_{\mathcal P}$,
    $S_\alpha=0$, and $S_{p_2}=0$.
    The selected implemented comparison also assumes
    $p_{\mathrm{ext}}\equiv0$; if $p_{\mathrm{ext}}$ varies with $z$, the
    wall/geometry source would need the additional term
    $-(a/\rho)\pd_z p_{\mathrm{ext}}$ in solver coordinates.
\end{definition}

\begin{proposition}[Rest-state compatibility of the selected wall source]
    \label{prop:selected-rest-state-compatibility}
    Suppose $p_{\mathrm{ext}}$ is constant in the gauge convention and the
    boundary data are compatible with $q=0$. Then
    $a(z)=R_0(z)^2$ and $q(z)=0$ are an exact steady state of the selected
    continuous source-balanced evolution law.
\end{proposition}

\begin{proof}
    The mass equation gives $\pd_z q=0$. In the momentum equation, all terms
    proportional to $q$ vanish, so
    $S_{\mathrm{fric}}=S_{\alpha}=S_{p_2}=0$. At $a=R_0^2$,
    \begin{flalign*}
        \quad
        \pd_z\Psi_h(R_0^2)
        &=
        \pd_z\left(\frac{K}{3\rho R_{\max}^2}R_0^3\right)
        =
        \frac{K}{\rho R_{\max}^2}R_0^2R_0'
        =
        S_{\mathrm{wall}}(R_0^2,z). &&
    \end{flalign*}
    Thus the elastic flux derivative is exactly canceled by the wall/geometry
    source, and $\pd_t q=0$.
\end{proof}

Continuous compatibility does not imply discrete preservation: the current
MUSCL/Rusanov realization exhibits the reported non-well-balanced
geometry-rest drift.

\subsection{Source-to-Implementation Differences}
\label{subsec:source-to-implementation-variant}

The selected implementation is a documented Canic-derived extended model, not
a literal reproduction of every source-paper parameter and normalization in
\cite{CanicEtAl2024Extended1DStenosis}. The report therefore uses the phrase
\emph{$R_{\max}$-normalized Canic-derived extended model} when source fidelity
matters. Table~\ref{tab:source-to-implementation-model} records the relevant
source-to-implementation choices. The diagnostic extended pressure correction is
\begin{flalign*}
    \quad
    p_{2,\mathrm g}(a,q,z,t)
    =
    g_{\mathcal P}\rho\nu_{\mathrm{eff}}^{\mathcal R}
    \frac{q}{a}\frac{R_0'(z)}{R_0(z)}. &&
\end{flalign*}
This $p_2$ term is a variable-radius viscous correction, not a replacement wall
law. In generalized-Newtonian descriptor runs, the implemented analytic $p_2$
derivative uses the local value of $\nu_{\mathrm{eff}}^{\mathcal R}$ and does
not differentiate that effective-viscosity factor; it is therefore a locally
frozen-viscosity computational closure for this term.

\begin{table}[!htb]
    \centering
    \scriptsize
    \caption{Source-to-implementation map for the selected 1D model.}
    \label{tab:source-to-implementation-model}
    \begin{tabularx}{\textwidth}{@{}
        >{\raggedright\arraybackslash}p{0.24\textwidth}
        >{\raggedright\arraybackslash}p{0.34\textwidth}
        >{\raggedright\arraybackslash}X@{}}
        \toprule
        \textbf{Feature} & \textbf{Canic et al. source form} & \textbf{Implemented report realization} \\
        \midrule
        Pressure denominator
            & Elastic pressure written with a local $R_0(z)^{-2}$ denominator.
            & Evolution wall potential and source use the constant
            $R_{\max}^{-2}$ normalization specified in
            Appendix~\ref{app:num-current-solver-operator}. \\
        Profile parameters
            & Source simulations state $\gamma=9$ and $\alpha=1.1$.
            & Principal runs use the parabolic profile
            $\alpha_{\mathcal P}=4/3$, $g_{\mathcal P}=4$; power-profile runs
            are sensitivity/diagnostic records. \\
        Friction coefficient
            & Profile-dependent wall/friction closure in the reduced model.
            & Implemented through the declared velocity-profile shear factor
            and rheology descriptor. \\
        $p_2$ differentiation
            & Variable-radius pressure correction differentiated in the model.
            & Locally frozen-viscosity realization for the $p_2$ derivative. \\
        External pressure
            & External pressure term may be included in the formulation.
            & $p_{\mathrm{ext}}\equiv0$ in the reported runs. \\
        Coordinates
            & Physical area and flow notation in the source paper.
            & Solver state is $(a,q)$ with
            $A_{\mathrm{phys}}=\pi a$, $Q_{\mathrm{phys}}=\pi q$. \\
        Boundary rule
            & Boundary conditions are model/problem dependent.
            & Fixed-area characteristic approximation applied with persisted
            subcritical-regime diagnostics. \\
        Discrete well balancing
            & Continuous rest equilibrium is a model-level property.
            & Rusanov/MUSCL is not exactly well balanced for spatially varying
            $R_0$; Chapter~\ref{sec:numerical-verification} reports
            zero-forcing rest-state drift. \\
        \bottomrule
    \end{tabularx}
\end{table}

\subsection{Boundary Approximation and Subcritical Condition}
\label{subsec:initial-boundary-conditions}

The comparison runs start from the geometry-rest state and then impose the
finite-volume boundary approximation stated below.

\begin{definition}[Single-vessel boundary data]
    \label{def:1d-single-vessel-boundary-data}
    Generic single-vessel models may prescribe inlet flow or velocity and an
    outlet pressure or terminal relation; these alternatives are model-family
    context rather than the reported boundary rule. The realized finite-volume
    run prescribes an inlet solver flow derived from a nominal reference-area
    mean velocity,
    \begin{flalign*}
        \quad
        q_{\mathrm{in}}&=R_0(0)^2\bar u_{\mathrm{in}}, &&
    \end{flalign*}
    solves the inlet ghost area from the approximate outgoing characteristic
    invariant, fixes the outlet ghost coordinate
    $a_{\mathrm{out}}=R_0(L)^2$, and obtains the outlet ghost flow from the
    corresponding approximate invariant. Hence the imposed inlet condition is
    a prescribed reference-area flow, not a fixed realized mean velocity
    $q_{\mathrm{in}}/a_{\mathrm{in}}$, and the outlet condition fixes the
    elastic reference pressure only for the selected wall component.
\end{definition}

\begin{definition}[Characteristic speeds for the reduced elastic subsystem]
    \label{def:1d-characteristic-speeds}
    Let $\lambda_\pm$ denote the two characteristic speeds of the
    positive-area elastic 1D subsystem obtained from the homogeneous
    $(A,Q)$ balance law. Their formula and assumptions are given in
    Appendix~\ref{app:num-elastic-potential-form}.
\end{definition}

\begin{assumption}[Subcritical boundary regime]
    \label{ass:1d-subcritical-boundary-regime}
    The baseline boundary treatment is used only in the subcritical regime
    $\lambda_-<0<\lambda_+$, where one characteristic enters through the inlet
    and one characteristic enters through the outlet.
\end{assumption}

Under this regime, a finite-volume scheme constrains the incoming information at
each boundary and leaves the outgoing state determined by the interior solve.
The finite-volume boundary realization in
Appendix~\ref{app:num-boundary-state-realization} uses an $\alpha=1$,
fixed-area characteristic approximation as the boundary invariant for the
variable-$\alpha_{\mathrm{eff}}$ solver. It is therefore an implemented boundary
approximation, not an exact invariant derivation for the full variable-radius
balance law.

\subsection{Principal MUSCL/Rusanov/SSPRK3 Realization}
\label{subsec:muscl-ssprk3-discretization}

Only the principal comparison method belongs in the main narrative. The reported
C23/C40 comparison uses the finite-volume realization specified in
Appendix~\ref{app:num-current-solver-operator}: the radius-squared state
$(a,q)$ is advanced on a one-dimensional grid with MUSCL reconstruction,
minmod-limited conservative states, Rusanov flux, source splitting for the
wall/geometry and variable-radius terms, and native SSPRK3 time integration.
The numerical comparison runs use $N=400$ cells on $0\le z\le6$ cm, a timestep
cap of $10^{-5}$ s, and a CFL cap of 0.45. The boundary-state approximation is
the fixed-area characteristic contract recorded in
Appendix~\ref{app:num-boundary-state-realization}. Other finite-volume, modal-DG,
native time-stepper, and SciML/OrdinaryDiffEq surfaces are implementation-check
or sensitivity context summarized in Table~\ref{tab:implemented-discretization-surface};
they are not the principal C23/C40 comparison method.

\subsection{Plane--Tetrahedron Observation Operator}
\label{subsec:cross-sectional-velocity-reconstruction}
\label{subsec:plane-tetrahedron-comparison-operator}

The 1D state directly supplies a section mean,
$\bar u_{1D}(z,t)=q(z,t)/a(z,t)$, and physical flow
$Q_{1D}(z,t)=\pi q(z,t)$. Principal radial profile plots use the
current-radius coordinate $r/\sqrt{a(z,t)}$ and the parabolic-profile
reconstruction
\begin{flalign*}
    \quad
    u_{1D}(r;z,t)
    &=
    2\,\frac{q(z,t)}{a(z,t)}
    \left(1-\left(\frac{r}{\sqrt{a(z,t)}}\right)^2\right) &&
\end{flalign*}
with the supplemental reference-radius coordinate $r/R_0(z)$ retained as an
operator-control output. That reconstruction is a profile diagnostic, not a
resolved secondary-flow, separation, skewness, or wall-traction model.

For a target plane $S_j=\Omega_{\mathrm{3D}}\cap\{z=z_j\}$, the default
\texttt{CrossSectionQuadratureOperator} intersects every tetrahedron with
$S_j$, linearly interpolates axial velocity on the cut edges, sorts the cut
polygon vertices in the plane, and triangulates the polygon by a centroid fan.
Let $\mathcal Q_j$ denote the resulting cut triangles. The reported 3D area,
flow, and mean velocity are
\begin{flalign*}
    \quad
    A_{3D,j}
    &=
    \sum_{\Delta\in\mathcal Q_j}|\Delta|,
    &
    Q_{3D,j}
    &=
    \sum_{\Delta\in\mathcal Q_j}|\Delta|\,\bar u_{z,\Delta},
    &
    \bar u_{3D,j}
    &=
    Q_{3D,j}/A_{3D,j}. &&
\end{flalign*}
For a cut triangle $\Delta$ whose vertices carry reconstructed axial
velocities $u_{z,1}$, $u_{z,2}$, and $u_{z,3}$,
\begin{flalign*}
    \quad
    \bar u_{z,\Delta}
    =
    \frac{u_{z,1}+u_{z,2}+u_{z,3}}{3}. &&
\end{flalign*}
The corresponding 1D quantities use physical flow, not the raw solver flow
coordinate:
\begin{flalign*}
    \quad
    Q_{1D,j}
    =
    \pi q(z_j,t),
    \qquad
    \bar u_{1D,j}
    =
    q(z_j,t)/a(z_j,t). &&
\end{flalign*}
Radial profile rows use the same cut triangles, binned by area-conservative
three-point triangle quadrature in either $r/\sqrt{a(z,t)}$ or $r/R_0(z)$.
The recorded controls include 10/20/40-bin sensitivity, area-weighted bin
means, within-bin azimuthal variance, and the current-to-reference area
mismatch. Empty bins are reported as invalid rows rather than zero-velocity
observations.

\subsection{Comparison Design, Discrepancy Metrics, and Provenance}
\label{subsec:numerical-experiments-metrics-provenance}

The resolved-velocity comparison evaluates the selected 1D solution against two
available resolved 3D velocity comparison datasets at nominally aligned
final-time samples: the 1D output is sampled at $1.0$ s and the XDMF velocity
field records $0.9995$ s. The cases are labelled C23 and C40 by stenosis
severity. The comparison is a diagnostic cross-model velocity comparison under
the declared observation operator. The C23/C40 section and radial quantities are
descriptive cross-model discrepancies under the declared operator, not
validation, accuracy, pressure-drop, FFR, physiological, clinical, predictive,
or causal evidence.

\begin{center}
\fbox{\begin{minipage}{0.90\linewidth}
\small
\textbf{Solver-coordinate convention.} In the comparison record, the stored
solver state \texttt{A} is $a=R^2$, not physical circular area. The physical
map is Definition~\ref{def:physical-to-solver-map}:
$A_{\mathrm{phys}}=\pi a$, $Q_{\mathrm{phys}}=\pi q$, and
$\bar u=q/a=Q_{\mathrm{phys}}/A_{\mathrm{phys}}$.
\end{minipage}}
\end{center}

The manuscript tables and figures use the tracked static report assets under
\path{report/assets/data/stenosis-comparison/}, interpreted under
Convention~\ref{conv:computed-realization-record}. The baseline regeneration
command, source-path convention, and file-hash manifest pointer are listed in
Appendix~\ref{app:code-and-ai-use}. Scratch sensitivity reruns live under
\path{julia/simulations/output/revision-*} or \path{tmp/revision-evidence/} and are
not the manuscript assets used for the current tables and figures. The raw
XDMF/HDF5 inputs remain local ignored files under
\path{julia/simulations/data/3d/canic_case3/}; the report archives the generated CSV
and PGFPlots-ready static data, not the raw fields.

\begin{table}[!htb]
    \centering
    \scriptsize
    \caption{Case matching for the resolved-velocity comparison. The case
    labels C23 and C40 are used in the main text and figures; the source ID
    records the upstream local data directory.}
    \label{tab:t1-case-matching}
    \begin{tabularx}{\textwidth}{@{}llXlrr@{}}
        \toprule
        \textbf{Label}
            & \textbf{Source ID}
            & \textbf{Velocity metadata}
            & \textbf{Operator}
            & \textbf{3D time (s)}
            & \textbf{Time offset (s)} \\
        \midrule
        C23 & 77 & \path{77/velocity.xdmf}
            & \texttt{CrossSectionQuadratureOperator}
            & 0.9994999999999453 & $5.0\times 10^{-4}$ \\
        C40 & 60 & \path{60/velocity.xdmf}
            & \texttt{CrossSectionQuadratureOperator}
            & 0.9994999999999453 & $5.0\times 10^{-4}$ \\
        \bottomrule
    \end{tabularx}
\end{table}

\begin{landscape}
\begin{table}[!htb]
    \centering
    \scriptsize
    \caption{Model-matching matrix for the 1D--3D velocity diagnostic. Status
    labels record what the archived report artifact supports, not what a future
    matched study might supply.}
    \label{tab:t1-model-matching-matrix}
    \begin{tabularx}{\linewidth}{@{}
        >{\raggedright\arraybackslash}p{0.13\linewidth}
        >{\raggedright\arraybackslash}p{0.23\linewidth}
        >{\raggedright\arraybackslash}p{0.23\linewidth}
        >{\raggedright\arraybackslash}p{0.11\linewidth}
        >{\raggedright\arraybackslash}X@{}}
        \toprule
        \textbf{Feature}
            & \textbf{1D model}
            & \textbf{3D comparison dataset}
            & \textbf{Status}
            & \textbf{Implication} \\
        \midrule
        Reference-profile geometry
            & Same analytic stenosis profile and declared severity.
            & Source mesh metadata is interpreted against the same nominal
            profile.
            & approximately matched
            & The intended stenosis shape is shared. \\
        Cut-operator area
            & Static reference area $A_{\mathrm{ref}}=\pi R_0^2$ is available.
            & Tetrahedral cuts audit against the same static reference area.
            & approximately matched
            & Cut areas show typical sub-percent agreement with localized
            multi-percent outliers. \\
        Current lumen area
            & Dynamic 1D area is $A_{1D}=\pi a(z,t)$.
            & Fixed 3D cut area is $A_{3D}(z)$.
            & algebraically reconstructed
            & The tracked section asset supports the reported flow/area split,
            but the split is bookkeeping rather than a wall-motion attribution. \\
        Wall model
            & Compliant reduced Canic--Koiter wall law.
            & Velocity and displacement XDMF files expose one local snapshot;
            wall-history and area-rate data are not persisted here.
            & unmatched
            & Wall effects and flow conservation cannot be isolated from the
            velocity mismatch. \\
        Density
            & $\rho$ recorded by the 1D parameters.
            & Not persisted in the tracked comparison tables.
            & unknown
            & Density-sensitive pressure or wave-speed claims are out of scope. \\
        Viscosity/rheology
            & Newtonian $\nu=0.04$ cm$^2$/s.
            & Rheology and viscosity metadata are not persisted in the
            report artifact.
            & unknown
            & Material-parameter matching is not established. \\
        Initial condition
            & Geometry-rest state, then forced by the inlet.
            & Prior transient history is not persisted.
            & unmatched
            & Final-time differences may include start-up history mismatch. \\
        Inlet condition
            & Reference-area flow
            $q_{\mathrm{in}}=R_0(0)^2\bar u_{\mathrm{in}}$ from nominal
            $\bar u_{\mathrm{in}}=22.5$ cm/s.
            & Boundary condition used for the 3D field is not persisted here.
            & unknown
            & Flow mismatch cannot be attributed uniquely to closure discrepancy. \\
        Outlet condition
            & Fixed-area characteristic approximation.
            & Outlet treatment is not persisted here.
            & unknown
            & Reflections and downstream constraints remain unquantified. \\
        Time / transient history
            & Evaluated at $T=1.0$ s from a start-up transient.
            & XDMF time is $0.9995$ s; earlier history is not recorded.
            & unknown
            & The final timestamps are close, but phase/history matching is not
            established. \\
        Numerical resolution
            & $N=400$ finite-volume cells.
            & Node-centered velocity fields with loaded node counts reported.
            & unresolved
            & Scratch reruns provide diagnostic resolution checks, but the
            manuscript tables remain single-realization descriptors unless a
            dedicated output-sensitivity study is completed. \\
        Velocity observable
            & Section mean $\bar u=q/a$ and physical flow $\pi q$.
            & Plane--tetrahedron cross-section quadrature of axial velocity.
            & matched
            & The compared observable is well defined. \\
        \bottomrule
    \end{tabularx}
\end{table}
\end{landscape}

\begin{table}[!htb]
    \centering
    \scriptsize
    \caption{Shared 1D realization and output-operator parameters for the
    $t=1.0$ comparison. The generated comparison records characteristic-speed
    diagnostics, while realized CFL, mass, and positivity diagnostics are
    exposed by the solver and verification records.}
    \label{tab:t1-parameter-matching}
    \begin{tabularx}{\textwidth}{@{}p{0.22\textwidth}X@{}}
        \toprule
        \textbf{Quantity} & \textbf{Recorded value} \\
        \midrule
        Forward model & \path{canic-extended-1d}. \\
        Wall model & $\mathcal W_{\mathrm{elastic1D}}$ with the selected
        Canic--Koiter pressure--area law
        \path{canic-koiter-thin-membrane}. \\
        Rheology & Newtonian rheology with $\nu=0.04$ cm$^2$/s. \\
        Velocity profile & Parabolic profile with
        $\alpha_{\mathcal P}=4/3$. \\
        Variable-radius terms & Canic effective-alpha and geometry-source
        corrections, including the locally frozen $p_2$ derivative in the
        extended member. \\
        Physical parameters & $\rho=1.055$ g/cm$^3$, $E=5.02\times10^6$
        dyn/cm$^2$, $h=0.06$ cm, $\sigma=0.5$,
        $K=4.016\times10^5$ dyn/cm, $R_{\max}=0.18$ cm, and
        $p_{\mathrm{ext}}\equiv0$. \\
        Discretization & $N=400$, $L=6$ cm, MUSCL finite volume with minmod
        limiter, native SSPRK3 time integration. \\
        Time controls & $T=1.0$ s from a start-up transient, timestep cap
        $10^{-5}$ s, CFL cap 0.45. \\
        Initial and boundary state & Geometry-rest initial state
        $a=R_0^2$, $q=0$; inlet $\bar u_{\mathrm{in}}=22.5$ cm/s; outlet
        $a_{\mathrm{out}}=R_0(L)^2$ with the implemented characteristic
        approximation in Appendix~\ref{app:num-boundary-state-realization}. \\
        Section targets & 200 equally spaced axial planes on $0\le z\le 6$
        cm. \\
        Radial targets & Slices $z=1.951$, 2.451, and 2.951 cm with principal
        20-bin plots in $r/\sqrt{a(z,t)}$ plus supplemental $r/R_0(z)$ and
        10/20/40-bin sensitivity rows. \\
        Supplemental sensitivity & Node-slab arithmetic means are stored only
        to \path{node-slab-sensitivity.csv} for half-widths 0.0075,
        0.0150753769, and 0.0301507538 cm. \\
        \bottomrule
    \end{tabularx}
\end{table}

Velocity and flow differences are reported as sample discrepancy measures over
valid section targets. With
$d_{u,j}=\bar u_{1D,j}-\bar u_{3D,j}$ and
$d_{Q,j}=Q_{1D,j}-Q_{3D,j}$, the retained section metrics are the signed mean
bias, mean absolute discrepancy, RMS discrepancy, maximum absolute discrepancy,
and relative RMS discrepancy,
\begin{flalign*}
    \quad
    \overline d_u
    &=
    \frac{1}{n}\sum_j d_{u,j},
    &
    D_{u,1}
    &=
    \frac{1}{n}\sum_j |d_{u,j}|,
    &
    D_{u,2}
    &=
    \left(\frac{1}{n}\sum_j d_{u,j}^2\right)^{1/2}, &&\\
    \quad
    D_{u,\infty}
    &=
    \max_j |d_{u,j}|,
    &
    D_{u,2}^{\mathrm{rel}}
    &=
    \frac{\left(\sum_j d_{u,j}^2\right)^{1/2}}
         {\left(\sum_j\bar u_{3D,j}^2\right)^{1/2}}, &&\\
    \quad
    \overline d_Q
    &=
    \frac{1}{n}\sum_j d_{Q,j},
    &
    D_{Q,1}
    &=
    \frac{1}{n}\sum_j |d_{Q,j}|,
    &
    D_{Q,2}
    &=
    \left(\frac{1}{n}\sum_j d_{Q,j}^2\right)^{1/2}. &&
\end{flalign*}
